\pdfbookmark[1]{Hassler. 1831 March 31.}{hassler1831}\par\Needspace{10\baselineskip}\begin{center}\rule{0.7\linewidth}{0.4pt}\par Göttingische gelehrte Anzeigen. 1831 March 31.\par\rule{0.7\linewidth}{0.4pt}\end{center}

\noindent\textit{Tabulae logarithmicae et trigonometricae notis septem decimalibus expressae.}
\textit{In forma minima. Purgatae ab erroribus praecedentium tabularum cura} F. K. \textsc{Hassler.} New York. G. C. and H. Carvill.

This transatlantic edition of the logarithm tables is distinguished by the exceptional neatness of its stereotype printing. Throughout, it gives to seven decimal places the logarithms of numbers up to 100000, and the logarithms of sines, tangents, cosines and cotangents for every second in the first degree, every ten seconds in the next two degrees, and every thirty seconds in all the remaining degrees of the quadrant. It also gives the natural sines and tangents throughout at intervals of thirty seconds. All this occupies 312 pages in small octavo or large duodecimo. Such compression could admittedly be achieved only with very small type, which, for all its beauty, may well be almost too small for everyday use by most eyes that are not nearsighted. Particular care seems to have been devoted to accuracy: at least, during a period in which we tried using the tables frequently, we encountered not a single printing error.

